% Partie : histoire
% Chapitre : biographies
% Section : Alfred Tarski

\documentclass[../../../include/open-logic-section]{subfiles}

\begin{document}

\olfileid[fr]{his}{bio}{tar}

\olsection{Alfred Tarski}

\olphoto{tarski-alfred}{Alfred Tarski}

Alfred Tarski naquit le 14 janvier 1901 à Varsovie, en Pologne, qui
faisait alors partie de l'Empire russe. Décrit comme «~napoléonien~»,
il était exubérant, bavard et passionné. Son énergie se manifestait
souvent pendant ses cours : un jour, en jetant une cigarette, il mit
le feu à une corbeille à papier et se vit interdire à l'avenir de
donner cours dans ce bâtiment.

Tarski avait soif de connaissances dès son plus jeune âge. Plus tard,
il raconterait à ses étudiants qu'il avait étudié la logique parce
que c'était la seule matière dans laquelle il avait obtenu la note~B.
Ses bulletins scolaires montrent pourtant qu'il avait obtenu des~A
partout, y compris en logique. Il étudia à l'université de Varsovie
de 1918 à 1924. Il avait d'abord l'intention d'étudier la biologie,
mais s'intéressa aux mathématiques, à la philosophie et à la logique,
car l'université était alors le centre de l'École de logique et de
philosophie de Varsovie. Il obtint son doctorat en 1924 sous la
direction de Stanis\l{}aw Le\'{s}niewski.

Avant d'émigrer aux États-Unis en 1939, Tarski réalisa certains de
ses travaux les plus importants alors qu'il enseignait dans le
secondaire à Varsovie. C'est à cette époque qu'il rédigea ses travaux
sur la conséquence logique et sur la définition de la vérité. En 1939, il se
trouvait aux États-Unis pour une tournée de conférences lorsque
l'Allemagne envahit la Pologne. En raison de ses origines juives, il
ne put rentrer. Sa femme et ses enfants restèrent en Pologne jusqu'à
la fin de la guerre, après quoi ils purent eux aussi émigrer aux
États-Unis. Tarski enseigna à Harvard, au College of the City of New
York et à l'Institute for Advanced Study de Princeton, puis à
l'université de Californie à Berkeley. Il y fonda le programme
interdisciplinaire de logique et de méthodologie des sciences.
Tarski mourut le 26 octobre 1983, à 82 ans.

\begin{reading}
Pour en savoir plus sur la vie de Tarski, voir la biographie
\emph{Alfred Tarski: Life and Logic} \citep{Feferman2004}.
Ses travaux fondamentaux sur la conséquence logique et la vérité
sont disponibles en anglais dans \citep{Tarski1983}. L'ensemble
de ses œuvres originales a été réuni en quatre volumes
\citep{Tarski1981}.
\end{reading}

\end{document}
